\subsection{The parity detector as a quotient of faithful quaternion homology}

\begin{result}{A.4.18}{Proved}{A literal characteristic homology target for the compact parity detector}
Use the exact groups of A.4.14 and the marked commutator
$c_\partial=[a,b]$.  Define
\[
 D_\partial=\langle\!\langle c_\partial^2\rangle\!\rangle_\Sigma,
 \quad L_c=L_8T'D_\partial,
 \quad R_c=L_c\cap\Sigma'\Sigma^{12},
 \quad\mathcal P_c=\Sigma/R_c.
                                                        \tag{QF.1}
\]
The actual orders are
\[
 |\Sigma/L_c|=8\cdot3^{14},\qquad
 |\mathcal P_c|=32\cdot3^{14}.                           \tag{QF.2}
\]
The restriction/deck detected modular image $J_c$ has
order $1152\cdot3^4$.  Exactly the same eighteen CM
residues as in (QH.5) lift to its centralizer.  Thus
this smaller target still excludes precisely the
opposite-parity classes among the six quotient cosets.

There is a specified characteristic quotient of the
original full profinite free surface group, with
literal surjection
\[
 \mathcal E^-\twoheadrightarrow\mathcal P_c,
 \qquad |\mathcal E^-|=288\cdot3^{36}.                  \tag{QF.3}
\]
It is obtained from the faithful quaternion summand
of $H_1(T,\mathbb F_3)$, joined to ordinary abelian
reduction modulo twelve.  The image of $T$ in
$\Sigma/L_c$ is $C_3^{12}$, and the exact source map
is a retained matrix
\[
 H_1(T,\mathbb F_3)^-\twoheadrightarrow C_3^{12}.
                                                        \tag{QF.4}
\]
In particular no nontrivial one-dimensional quaternion
character block or ordinary compact-period coordinate
is silently included in (QF.4).  The signed period
information is the separate abelian factor in (QF.3).

More precisely, let $H=C_3^2$ be the specified translation
action, and let $\mathfrak m_H$ be the augmentation ideal
of $\mathbb F_3[H]$.  For $A=H_1(T,\mathbb F_3)$,
the actual map $\Theta:A\to\mathbb F_3^{14}$ below has
the exact kernel
\[
 \ker\Theta=A^+\oplus\mathfrak m_H^2A^-.
                                                       \tag{QF.K1}
\]
On the primitive eighteen-dimensional quaternion
multiplicity space $M$, the exact quotient is
$M/\mathfrak m_H^2M$, of dimension six.  Every linear
normalizer of the given $H$ action preserves this
particular quotient kernel.

This theorem supplies the common finite quotient and
the CM centralizer test.  It does not identify the
normalizer of the finite translation matrices in the
orthogonal image with the image of the full geometric
deck normalizer; the joint finite comparison still
retains that requirement.
\end{result}

\begin{proof}
The $36$ peripheral loops of $T$ are precisely conjugates
in $\Sigma$ of $c_\partial^2$, with the usual choice of
one conjugating path for each boundary component.
Indeed $\Sigma/T=Q_8\times C_3^2$ has order $72$,
and the original peripheral loop maps to the central
involution $(-1,0)$.  Its inertia therefore has order
two and gives $72/2=36$ lifts.  This proves that
$D_\partial$ in (QF.1) is exactly the normal closure of
those boundary loops, with the specified based map;
it is not a different compactification.

Begin with the explicit $24$-dimensional Lie group
of (QH.8), using the class-two multiplication with
$1/2=2\in\mathbb F_3$.  Rewrite the actual word
$c_\partial^2$ in the nine Schreier generators of
$N_8$, and collect it in these coordinates.  The
program forms the linear span of this vector, its
images under the four actual automorphisms of (QS.10),
and all its brackets with the nine abelian basis
vectors, repeating until stable.  This is exactly
the normal Lie ideal generated by that word.
The two deck substitutions together with conjugation
by $N_8$ supply all conjugations by $\Sigma$.
The two modular substitutions preserve the peripheral
normal closure, so adjoining them does not enlarge it
beyond that normal closure.  Class-two multiplication
and its inverse identify this ideal with the normal
subgroup it generates: brackets are already retained,
and division by two is allowed in characteristic
three.  All these claims are checked by the retained
full matrices and projection identities.

The resulting ideal has dimension ten.  The quotient
has abelian dimension six and central dimension eight.
The program records an actual $14\times24$ projection,
a section as vector spaces, and the four induced
$14\times14$ automorphism matrices; their intertwining
identities hold on all $24$ source coordinates.
This proves $|N_8/L_c|=3^{14}$ and the first equality
in (QF.2).  Every added peripheral word has zero
ordinary exponent sums, so the exponent lattice of
$L_c$ is still $6\mathbb Z^2$.  The ordinary abelian
join adds four, proving the second equality.
The normal subgroup is invariant under the actual
modular automorphisms, which preserve the peripheral
commutator.  Full characteristic invariance will be
supplied by the different, explicitly specified
source $\mathcal E^-$ below.

\textbf{The exact compact centralizer census.}
The six abelian and eight central quotient coordinates
give a literal bracket map
$\bigwedge^2\mathbb F_3^6\to\mathbb F_3^8$.
The program retains this map and verifies its rank
is eight.  The inner derivations have dimension six,
so the outer central-coordinate space has dimension
$8\cdot6-6=42$.  The four actual substitutions,
with their original signed mod-twelve matrices,
give the same $9216$-element joint linear base as
before.  The complete $36864$ Schreier rows have
rank four, and their span is stable under all four
generators.  The exact image therefore has order
$1152\cdot3^4$ after quotienting by the eight actual
deck elements.

For every one of the $36$ CM-commuting base residues,
the program records the actual source word and deck
correction and solves its four-dimensional version
of (QS.15).  The coefficient rank is two.  The six
quotient-coset obstruction vectors are again those
of (QH.11).  There are eighteen surviving residues,
each with nine lifts; the excluded classes include
$5I$.  This proves the stated centralizer assertion
for $J_c$, by the same exact affine criterion.

\textbf{The actual faithful homology map.}
The original exponent map on the nine Schreier words
induces a rank-two map on the compact quotient's
fourteen coordinates.  Its kernel is the exact
image of $T$, of dimension twelve.  To verify both
the source and the factorization, the program takes
each of the $73$ retained basis words of $T$ in
$\Pi_{1,9}$, substitutes its actual word in $\Sigma$,
rewrites it in $N_8$, and applies the displayed
$14$-coordinate projection.  This gives a retained
matrix $\Theta$ of size $14\times73$, with rank twelve.
It satisfies
\[
 \Theta(I+D_{-1})=0,\qquad
 \ell_c\Theta=0,                                     \tag{QF.5}
\]
where $D_{-1}$ is the actual central quaternion deck
matrix from A.4.15 and $\ell_c$ is the specified
rank-two period map.  The first identity says exactly
that the map kills the positive quaternion summands.
Equivalently, the actual central involution acts
as $-I$ on the twelve-dimensional target.
For instance its based lift $a^6$ has quaternion
label $-1$ and trivial three-division label, and its
sixfold conjugation matrix gives this same action.
Thus (QF.5), with rank twelve, proves the surjection
(QF.4), including its based source map.

\textbf{The exact augmentation-square kernel.}
The independent verifier
\path{work/ghost_code/check_main_faithful_detector.py}
rebuilds quaternion multiplication from
$i^2=j^2=-1$, $ji=-ij$, and independently reconstructs
both the nine- and seventy-three-word Schreier bases.
It performs every substitution defining $\Theta$,
using the class-two law
$(v,z)(w,t)=(v+w,z+t+\tfrac12v\wedge w)$,
then the actual projections of dimensions
$36\to15$ and $24\to14$.
The resulting $14\times73$ matrix has rank twelve.
Every bracket between its seventy-three generator
images is zero.  The target three-group has exponent
three, so these identities establish the factorization
through $T/(T'T^3)$ as a group homomorphism.
The same verifier rewrites all thirty-six actual
peripheral lifts.  Their homology vectors span
dimension thirty-five, and all map to zero.
It also checks (QF.5) and
$\Theta\,\tfrac12(I-D_{-1})=\Theta$ entry by entry.

Let $h_1,h_2$ denote the actual translation matrices
on $A$, and let $C$ be the retained $73\times36$
basis matrix of $A^-$.  The group algebra is
\[
 \mathbb F_3[H]\cong
 \mathbb F_3[X,Y]/(X^3,Y^3),\qquad
 X=h_1-I,\quad Y=h_2-I.
                                                       \tag{QF.K2}
\]
In particular $\mathfrak m_H^2$ is generated by
$X^2,XY,Y^2$.
The verifier constructs the actual image matrix
\[
 \mathcal J=
 \bigl[(h_1-I)^2C\ \ (h_1-I)(h_2-I)C\ \ (h_2-I)^2C\bigr].
                                                       \tag{QF.K3}
\]
It checks $\Theta\mathcal J=0$,
$\operatorname{rank}\mathcal J=24$, and
$\operatorname{rank}(\Theta C)=12$.
The first equality gives inclusion of the entire
augmentation-square image in the kernel on $A^-$;
the two ranks give equality, because $\dim A^-=36$.
Together with (QF.5), this proves both inclusions
in (QF.K1).

For the retained primitive multiplicity basis
$C_B$ of size $73\times18$, direct substitution
gives $\operatorname{rank}(\Theta C_B)=6$.
On its eighteen-dimensional domain the three
matrices $(H_1-I)^2,(H_1-I)(H_2-I),(H_2-I)^2$
have joint rank twelve and are killed by
$\Theta C_B$.  Their image is
$\mathfrak m_H^2M$ by (QF.K2), so the same
dimension argument proves that this is the exact
kernel and proves the asserted six-dimensional
quotient with its actual map.

Finally any linear normalizer $g$ of the specified
$H$ action satisfies, for $h,k\in H$ and $v\in M$,
\[
 g(h-I)(k-I)v
 =(ghg^{-1}-I)(gkg^{-1}-I)gv.
                                                       \tag{QF.K4}
\]
This takes $\mathfrak m_H^2M$ onto itself,
using the same equation for $g^{-1}$.
It therefore induces an actual automorphism
of $M/\mathfrak m_H^2M$.
For a map carrying $H$ to a specified $H'$,
the identical equation gives precisely the
isomorphism $M/\mathfrak m_H^2M\to
M/\mathfrak m_{H'}^2M$.
These are statements about the displayed module
quotient; the theorem's joint signed-period
normalizer requirement remains in place.
The independent NPZ record retains $\Theta$,
all peripheral rows, the faithful projector,
and both exact augmentation-square image matrices.

\textbf{The characteristic common source.}
The group $\Sigma/T=Q_8\times C_3^2$ has a unique
involution $z=(-1,0)$.  Its action on
$A=T/(T'T^3)$ is independent of the chosen lift,
since conjugation by $T$ is trivial on this module.
Let $A^+$ be its positive eigenspace, and define
\[
 K^- =\ker\bigl(T\to A/A^+\bigr),\qquad
 \mathcal E^-=
       \Sigma/(K^-\cap\Sigma'\Sigma^{12}).              \tag{QF.6}
\]
The subgroup $T$ is characteristic in $\Sigma$,
the element $z$ is characteristic in $\Sigma/T$,
and the positive eigenspace is therefore invariant
under every induced automorphism.  Consequently
$K^-$ and its displayed intersection are
characteristic, also for continuous automorphisms
of the full profinite free group.
The eigenspace calculation gives
$|T/K^-|=3^{36}$, so
$|\Sigma/K^-|=72\cdot3^{36}$.

For completeness the ordinary exponent lattice of
$K^-$ is precisely $6\mathbb Z^2$.  Containment
follows from $K^-\subseteq T$.  The words
$a^6c_\partial$ and $b^6c_\partial$ lie in $T$ and
have exponent vectors $(6,0)$ and $(0,6)$.
For either such word $w$, and any lift $\tilde z$
of $z$, both
$w\tilde z w\tilde z^{-1}$ and $w^3$ belong to
$K^-$: the first has homology class $(I+z)[w]$,
and the second dies in $A$.
Their exponent vectors are respectively twice and
three times that of $w$.  Integer combinations
therefore give each of $(6,0),(0,6)$ in the image
of $K^-$.  Joining ordinary abelian reduction modulo
twelve adds exactly four, proving (QF.3)'s order.

Finally (QF.5) proves $K^-\subseteq L_c$: the map
$T\to N_8/L_c$ factors through $A$ and kills $A^+$.
Intersecting this inclusion with
$\Sigma'\Sigma^{12}$ gives the actual surjection
(QF.3).  This establishes the promised common
characteristic source, not merely equality of
some module dimensions.

The complete sources are
\path{work/tower_code/modular_quaternion_compact_setup.py},
\path{work/tower_code/modular_quaternion_compact_census.py},
\path{work/tower_code/modular_quaternion_compact_faithful.py},
and their matching output, JSON and NPZ files.
They retain the ideal, all quotient projections,
the literal $14\times73$ map, the central involution,
and every affine centralizer equation.
\end{proof}
